\subsection{EXAMPLE 5.5: FRM EXAM 2004---QUESTION 39 (portfolio volatility)}

\textbf{Enunciado}

\emph{Consider a portfolio with 40\% invested in asset X and 60\% invested in
asset Y. The mean and variance of return on X are 0 and 25
respectively. The mean and variance of return on Y are 1 and 121
respectively. The correlation coefficient between X and Y is 0.3. What is
the nearest value for portfolio volatility?}
\begin{enumerate}[label=\alph*.]
\item 9.51
\item 8.60
\item 13.38
\item \hl{7.45}
\end{enumerate}

\subsubsection*{Solución}

Sea \(R_P\) el rendimiento del portafolio. Como el 40\% está invertido en \(X\) y el 60\%
en \(Y\), se tiene
\[
R_P = 0.4\,R_X + 0.6\,R_Y.
\]

Además, se nos dan las varianzas, desviaciones estándar y correlación de los rendimientos de \(X\) e \(Y\):
\[
\Var(R_X)=25,\qquad \Var(R_Y)=121,
\]
\[
\sigma_X=\sqrt{25}=5,\qquad \sigma_Y=\sqrt{121}=11.
\]\[
\rho_{XY}=0.3.
\]

La fórmula de la varianza de un portafolio de dos activos para constantes \(a,b\) es:
\[
\Var(aA+bB)=a^2\Var(A)+b^2\Var(B)+2ab\,\mathrm{Cov}(A,B).
\]
\[
\Var(R_P)
=
w_X^2\sigma_X^2+w_Y^2\sigma_Y^2+2w_Xw_Y\rho_{XY}\sigma_X\sigma_Y.
\]
\[
\sigma_P^2
=
(0.4)^2(5)^2+(0.6)^2(11)^2+2(0.4)(0.6)(0.3)(5)(11).
\]
\[
\sigma_P^2=4+43.56+7.92=55.48.
\]

Finalmente, la volatilidad del portafolio es
\[
\sigma_P=\sqrt{55.48}\approx 7.45.
\]

\paragraph{Observación.}
Si la correlación hubiera sido \(\rho=1\), entonces la volatilidad sería
\[
0.4(5)+0.6(11)=8.6,
\]
pero como \(\rho=0.3<1\), hay beneficio de diversificación y la volatilidad
resulta menor.








